\documentclass[10pt,a4paper]{amsart}
\usepackage{amsmath,amssymb,amsthm,graphicx}
% For commutative diagrams we use Xy-pic
\usepackage[all]{xy}

% We use 2cell for 2-commutative diagrams.
\xyoption{2cell}
% Theorem environments.
%
\theoremstyle{plain}
\newtheorem{theorem}[subsection]{Theorem}
\newtheorem{proposition}[subsection]{Proposition}
\newtheorem{lemma}[subsection]{Lemma}

\theoremstyle{definition}
\newtheorem{definition}[subsection]{Definition}
\newtheorem{example}[subsection]{Example}
\newtheorem{exercise}[subsection]{Exercise}
\newtheorem{situation}[subsection]{Situation}

\theoremstyle{remark}
\newtheorem{remark}[subsection]{Remark}
\newtheorem{remarks}[subsection]{Remarks}

\numberwithin{equation}{subsection}

% Macros
%
\def\lim{\mathop{\mathrm{lim}}\nolimits}
\def\colim{\mathop{\mathrm{colim}}\nolimits}
\def\Spec{\mathop{\mathrm{Spec}}}
\def\Hom{\mathop{\mathrm{Hom}}\nolimits}
\def\Ext{\mathop{\mathrm{Ext}}\nolimits}
\def\SheafHom{\mathop{\mathcal{H}\!\mathit{om}}\nolimits}
\def\SheafExt{\mathop{\mathcal{E}\!\mathit{xt}}\nolimits}
\def\Sch{\mathit{Sch}}
\def\Mor{\mathop{\mathrm{Mor}}\nolimits}
\def\Ob{\mathop{\mathrm{Ob}}\nolimits}
\def\Sh{\mathop{\mathit{Sh}}\nolimits}
\def\NL{\mathop{N\!L}\nolimits}
\def\CH{\mathop{\mathrm{CH}}\nolimits}
\def\proetale{{pro\text{-}\acute{e}tale}}
\def\etale{{\acute{e}tale}}
\def\QCoh{\mathit{QCoh}}
\def\Ker{\mathop{\mathrm{Ker}}}
\def\Im{\mathop{\mathrm{Im}}}
\def\Coker{\mathop{\mathrm{Coker}}}
\def\Coim{\mathop{\mathrm{Coim}}}

% Boxtimes
%
\DeclareMathSymbol{\boxtimes}{\mathbin}{AMSa}{"02}

%
% Macros for moduli stacks/spaces
%
\def\QCohstack{\mathcal{QC}\!\mathit{oh}}
\def\Cohstack{\mathcal{C}\!\mathit{oh}}
\def\Spacesstack{\mathcal{S}\!\mathit{paces}}
\def\Quotfunctor{\mathrm{Quot}}
\def\Hilbfunctor{\mathrm{Hilb}}
\def\Curvesstack{\mathcal{C}\!\mathit{urves}}
\def\Polarizedstack{\mathcal{P}\!\mathit{olarized}}
\def\Complexesstack{\mathcal{C}\!\mathit{omplexes}}
% \Pic is the operator that assigns to X its picard group, usage \Pic(X)
% \Picardstack_{X/B} denotes the Picard stack of X over B
% \Picardfunctor_{X/B} denotes the Picard functor of X over B
\def\Pic{\mathop{\mathrm{Pic}}\nolimits}
\def\Picardstack{\mathcal{P}\!\mathit{ic}}
\def\Picardfunctor{\mathrm{Pic}}
\def\Deformationcategory{\mathcal{D}\!\mathit{ef}}
\usepackage{cleveref}
\crefname{originallemma}{original lemma}{original lemmas}
\Crefname{originallemma}{Original lemma}{Original lemmas}
\makeatletter
\newlabel{algebra-lemma-lift-etale}{{04D1}{source}{}{sourcepointer-1}{}}
\newlabel{algebra-lemma-lift-etale@cref}{{[originallemma][0][]04D1}{[1][0][]source}{}{}{}}
\makeatother
\crefname{originalreference}{original reference}{original references}
\Crefname{originalreference}{Original reference}{Original references}
\makeatletter
\newlabel{lemma-normal-scheme-with-alg-closed-function-field}{{09Z9}{source}{}{sourcepointer-2}{}}
\newlabel{lemma-normal-scheme-with-alg-closed-function-field@cref}{{[originalreference][0][]09Z9}{[1][0][]source}{}{}{}}
\makeatother
\title{Etale covers of a closed subscheme of an affine normal scheme with separably closed function field split into open pieces, and finite surjective ones have sections}
\author{The Stacks Project Authors}
\date{stacks-project:etale-cohomology.tex:local-snapshot-20261011; GNU FDL 1.2 or later}
\begin{document}
\maketitle
\noindent\textit{Editorial note.} All delivered mathematical text is the author's own, copied from the pinned Stacks Project chapter file etale-cohomology.tex; no statement or proof step was rewritten, reordered or removed.\par\medskip
\noindent\textit{Editorial note.} The counted proof is the author's complete argument: the proof is written out in full: the lift to an affine etale cover of X, the application of the structure lemma, the decomposition into open immersions that are also closed, and the displayed open-and-closed decomposition producing a section.\par\medskip
\noindent\textit{Editorial note.} The label lemma-closed-of-affine-normal-scheme-with-alg-closed-function-field names the counted statement, and the proof environment immediately following it in the source is delivered verbatim as its proof.\par\medskip
\noindent\textit{Editorial note.} The counted statement is carried by the author's own lemma environment in the chapter source.\par\medskip
\noindent\textit{Editorial note.} The statement and proof are copied from the pinned chapter file; the two cited inputs are kept as external pointers with their canonical Stacks tags.\par\medskip
\noindent\textit{Allowed standard prerequisites (editorial).}\begin{itemize}
\item Lifting an etale morphism from a closed subscheme to the ambient scheme. Algebra, the lemma allowing an etale morphism V -> Z from a closed subscheme to be lifted to an affine scheme U etale over X. Supplies the lift U -> X to which the normal-scheme structure statement is applied.
\item Separated etale morphisms over a normal scheme with separably closed function field. Same chapter, the lemma showing that on such a scheme a separated etale morphism is a disjoint union of open immersions. Supplies the first statement, applied to U -> X.
\end{itemize}\medskip
\ifdefined\hypertarget\hypertarget{sourcepointer-1}{}\fi
\noindent\textit{Source reference.} Original lemma 04D1: sources/stacks-chow/tags.txt, source lines 5711--5711. Prerequisite pointer invoked at proof line(s) 14399; it names a standard fact of a sibling chapter that the counted argument uses, and every step of the counted proof built on it is retained verbatim in the delivered body.\par\medskip
\ifdefined\hypertarget\hypertarget{sourcepointer-2}{}\fi
\noindent\textit{Source reference.} Original reference 09Z9: sources/stacks-etale-cohomology/pinned-original-source.tex, source lines 14312--14312. Prerequisite pointer invoked at proof line(s) 14401; it names a standard fact of the same chapter that the counted argument uses, and every step of the counted proof built on it is retained verbatim in the delivered body.\par\medskip
\begin{lemma}
\label{lemma-closed-of-affine-normal-scheme-with-alg-closed-function-field}
Let $X$ be an affine integral normal scheme with separably closed
function field. Let $Z \subset X$ be a closed subscheme. Let
$V \to Z$ be an \'etale morphism with $V$ affine. Then $V$ is a finite
disjoint union of open subschemes of $Z$. If $V \to Z$ is
surjective and finite \'etale, then $V \to Z$ has a section.
\end{lemma}

\begin{proof}
By Algebra, Lemma \ref{algebra-lemma-lift-etale}
we can lift $V$ to an affine scheme $U$ \'etale over $X$.
Apply Lemma \ref{lemma-normal-scheme-with-alg-closed-function-field}
to $U \to X$ to get the first statement.

\medskip\noindent
The final statement is a consequence of the first.
Let $V = \coprod_{i = 1, \ldots, n} V_i$ be a finite
decomposition into open and
closed subschemes with $V_i \to Z$ an open immersion.
As $V \to Z$ is finite we see that $V_i \to Z$ is also closed.
Let $U_i \subset Z$ be the image. Then we have a decomposition
into open and closed subschemes
$$
Z =
\coprod\nolimits_{(A, B)}
\bigcap\nolimits_{i \in A} U_i \cap
\bigcap\nolimits_{i \in B} U_i^c
$$
where the disjoint union is over $\{1, \ldots, n\} = A \amalg B$
where $A$ has at least one element.
Each of the strata is contained in a single $U_i$ and
we find our section.
\end{proof}

\end{document}
